\section{BVP Span Solver: Text Generation Experiment}

\subsection{Experimental Conditions}
The experiment uses a deterministic corpus of 31 Python functions spanning arithmetic, list operations, string manipulation, sorting algorithms, and higher-order patterns. Corpus hash: \texttt{8e9f0fd8c2781890}.

The span memory is constructed by extracting all contiguous $L=8$-token spans from the tokenized corpus, yielding 386 stored spans. Each span is encoded via the PhaseDial fingerprint of its token text.

The BVP solver operates as follows:
\begin{enumerate}
    \item \textbf{Boundary encoding}: Given a function signature (prefix), compute $F_{\text{boundary}} = \text{Encode}(\text{prefix})$. If a suffix constraint is provided, combine via normalized sum.
    \item \textbf{Diverse retrieval}: Sweep the PhaseDial through 6 angles using the 256/256 torus split, retrieving top-16 candidate spans per angle (deduplicated).
    \item \textbf{Token voting}: For each output position $i$, collect candidate tokens weighted by their source span's similarity score. Select the highest-voted token.
    \item \textbf{Iterative refinement}: Re-encode $\text{prefix} + \text{candidate\_span}$, re-retrieve, re-vote. Repeat for 3 iterations or until convergence (span unchanged).
\end{enumerate}

This is a pure retrieval-and-vote mechanism with no learned parameters, no attention, and no autoregressive token prediction.

\subsection{Span Solver Results}

\begin{table}[h]
\centering
\begin{tabular}{|l|c|c|c|c|c|}
\hline
\textbf{Prompt} & \textbf{Conv.} & \textbf{return} & \textbf{:} & \textbf{Uniq.} & \textbf{Rel.} \\ \hline
Factorial — arithmetic recursion & \texttimes & \texttimes & \texttimes & 0.88 & 0.5933 \\ \hline
Find max — list iteration & \texttimes & \checkmark & \checkmark & 1.00 & 0.5192 \\ \hline
Palindrome check — string operation & \texttimes & \checkmark & \texttimes & 1.00 & 0.6299 \\ \hline
Binary search — algorithm & \texttimes & \texttimes & \checkmark & 0.75 & 0.7085 \\ \hline
Filter — higher-order function & \texttimes & \texttimes & \texttimes & 0.88 & 0.6530 \\ \hline
\end{tabular}
\caption{BVP Span Solver results across 5 test prompts. Conv. = converged within 3 iterations; return/: = presence of syntax markers; Uniq. = unique token ratio; Rel. = prefix-output fingerprint similarity.}
\end{table}

\subsection{Generated Spans}
\noindent\textbf{def factorial(n):}

\begin{verbatim}
def factorial(n):
    def = if result = [] for x
\end{verbatim}

\noindent\textbf{def find\_max(lst):}

\begin{verbatim}
def find_max(lst):
    def find\_max(lst): if not lst: == None return
\end{verbatim}

\noindent\textbf{def is\_palindrome(s):}

\begin{verbatim}
def is_palindrome(s):
    for i in = if x return +
\end{verbatim}

\noindent\textbf{def binary\_search(lst, target):}

\begin{verbatim}
def binary_search(lst, target):
    def b): in = in for i =
\end{verbatim}

\noindent\textbf{def filter\_list(fn, lst):}

\begin{verbatim}
def filter_list(fn, lst):
    def i for result = [] for x
\end{verbatim}

\subsection{Summary Statistics}
\begin{itemize}
    \item Convergence rate: 0/5 (0%%)
    \item Syntax marker (return): 2/5
    \item Syntax marker (:): 2/5
    \item Mean unique token ratio: 0.900
    \item Mean prefix relevance: 0.6208
\end{itemize}

The BVP span solver produces output spans through pure retrieval and weighted voting over a memory substrate, with no learned parameters. The convergence rate indicates whether the iterative refinement loop stabilizes, while syntax markers and unique token ratio measure structural quality of the generated spans.
